\begin{preface}{TÓM TẮT}
Sự phát triển của học máy và điện toán biên đã cách mạng hóa các ứng dụng theo dõi chuyển động, cho phép các hệ thống thời gian thực, tiết kiệm năng lượng và chính xác cao cho nhiều trường hợp sử dụng trong y tế, phân tích thể thao và tương tác người-máy. Tuy nhiên, việc phát triển các mô hình AI cho theo dõi chuyển động vẫn còn phức tạp, đòi hỏi kiến thức chuyên môn về tiền xử lý dữ liệu, trích xuất đặc trưng, huấn luyện mô hình và triển khai trên thiết bị biên. Bài báo này trình bày một framework toàn diện giúp đơn giản hóa quy trình end-to-end để tạo ra các mô hình AI được tùy chỉnh cho theo dõi chuyển động con người trên các thiết bị biên có tài nguyên hạn chế.

Framework giới thiệu một phương pháp tiền xử lý tương tác mới với tính năng lựa chọn cửa sổ thời gian kéo thả, cho phép người dùng chú thích và phân đoạn dữ liệu chuyển động một cách hiệu quả mà không cần kiến thức lập trình sâu rộng. Hệ thống hỗ trợ nhiều thuật toán học máy bao gồm Random Forest, Support Vector Machines và Neural Networks, với khả năng tự động sinh mã cho các nền tảng biên khác nhau bao gồm vi điều khiển Arduino, ARM Cortex-M và Seeed XIAO. Việc sinh mã nhận biết tối ưu hóa của chúng tôi xem xét nhiều chiến lược triển khai—độ chính xác, tốc độ, tiêu thụ điện năng và các phương pháp cân bằng—điều chỉnh mã được sinh ra theo yêu cầu ứng dụng cụ thể.

Chúng tôi xác thực framework bằng cách sử dụng bộ dữ liệu Nhận dạng Hoạt động Con người (HAR) được thu thập ở 100Hz từ cảm biến IMU 6 trục (LSM6DS3) trên nền tảng Seeed XIAO nRF52840 Sense. Trong nghiên cứu bằng chứng khái niệm này sử dụng dữ liệu đơn đối tượng trên năm lớp hoạt động, kết quả thực nghiệm cho thấy các mô hình đã triển khai đạt độ chính xác 94.5\% cho Random Forest và 96.2\% cho Neural Networks, với thời gian suy luận lần lượt là 15ms và 12ms, trong khi tiêu thụ ít hơn 10mW trong quá trình hoạt động. Phân tích so sánh với các nền tảng thương mại (Edge Impulse, SensiML) cho thấy framework mã nguồn mở của chúng tôi cung cấp tính linh hoạt tiền xử lý vượt trội và chất lượng sinh mã tốt hơn với chi phí bằng không.

Framework giải quyết các khoảng trống quan trọng trong các giải pháp hiện tại bằng cách cung cấp: (1) quy trình làm việc dựa trên GUI trực quan dễ tiếp cận cho người không chuyên, (2) kiểm soát hoàn toàn đường ống tiền xử lý với phản hồi trực quan, (3) triển khai không phụ thuộc phần cứng với các tối ưu hóa cụ thể cho từng nền tảng, và (4) kiến trúc có thể mở rộng hỗ trợ tích hợp các thuật toán và nền tảng bổ sung trong tương lai. Nghiên cứu này đóng góp vào việc dân chủ hóa phát triển AI biên cho các ứng dụng theo dõi chuyển động, với tác động tiềm năng trong phục hồi chức năng y tế, giám sát hiệu suất thể thao và công nghệ hỗ trợ.
\end{preface}
